\documentclass[11pt,a4paper]{article}
\usepackage[utf8]{inputenc}
\usepackage[T1]{fontenc}
\usepackage[english,russian]{babel}
\usepackage{amsmath,amssymb,amsthm}
\usepackage{geometry}
\geometry{margin=2.2cm}

\newtheorem{theorem}{Theorem}
\newtheorem{lemma}[theorem]{Lemma}
\newtheorem{corollary}[theorem]{Corollary}
\newtheorem{definition}[theorem]{Definition}
\newtheorem{remark}[theorem]{Remark}
\newtheorem{axiom}{Axiom}

\newcommand{\subfact}{\mathord{!}\,}
\newcommand{\S}{\mathcal{S}}

\begin{document}

\title{Вероятностно-комбинаторная логика времени:\\
теорема о маятнике со стрелой (v7.1)}
\author{Алекс Мохов}
\date{15.11.2026}

\maketitle

\begin{abstract}
Формулируется центральная теорема Вероятностно-комбинаторной семизначной логики
времени (7VL): время как период/апериод смены ролей $L_4 \leftrightarrow L_3$
выводится из пяти сцепленных утверждений, все числовые константы которых
($2$, $3$, $4$, $7$, $11$, $12$, $13$, $15$, $16$) не содержат свободных
параметров. Новая часть --- п.~(v): стрела времени. Доказано, что марковские
каналы расширения $3 \to 4$ и стягивания $4 \to 3$ имеют вероятности удержания
инварианта $\frac{13}{16}$ и $\frac{13}{15}$ соответственно, что нарушает
принцип детального баланса с точным коэффициентом $\chi = \frac{16}{15}$.
Стационарное распределение двухтактной цепи равно
$\left(\frac{16}{31},\,\frac{15}{31}\right)$, а квант асимметрии потерь равен
$\frac{13}{240}$ --- числитель совпадает с модулем сброса PG$(2,3)$.
\end{abstract}

\section{Аксиоматика и обозначения}

\begin{axiom}[Регистр состояния]
Элементарный тактовый регистр есть ортогональная сумма
$R_7 = L_4 \oplus L_3$, $\dim R_7 = 4 + 3 = 7$, где $L_4$ --- тетрарная логика
данных, $L_3$ --- тернарная логика контроля.
\end{axiom}

\begin{definition}[Время]
Время $\tau$ --- длительность удержания текущего распределения ролей
$L_4 \leftrightarrow L_3$ до акта их взаимной инверсии. Штатный ход ---
период $T_0$; превышение порога ошибок --- апериод (скачок-сброс).
\end{definition}

Для $\S_n$: числа встреч $D_{n,j} = \binom{n}{j}\subfact{(n-j)}$;
$\Delta_n(k) = \sum_{j \ge k} D_{n,j}$; $\Delta_n \equiv \Delta_n(1) = n! - \subfact{n}$.

\section{Теорема о маятнике со стрелой}

\begin{theorem}[Вероятностно-комбинаторная логика времени]\label{th:main}
Единственной самосогласованной дискретной системой с коррекцией ошибок, в которой
пространство данных и пространство контроля симметрично обмениваются ролями и
которая порождает ориентированное (необратимое) собственное время, является
замкнутый двухтактный маятник $3 \leftrightarrow 4$ со следующими пятью
сцепленными характеристиками:
\begin{enumerate}
\item[(i)] \textbf{Регистрация.} Двоичный алфавит $\{0,1\}$ задается
$\subfact{3} = 2$ --- единственным простым (и единственным четным простым) среди
всех субфакториалов.
\item[(ii)] \textbf{Баланс.} Уравнение насыщения каналов
$\Delta_{n+1} - \Delta_n = n(n+1) - 1$ имеет единственный натуральный корень
$n = 3$ при $n \ge 3$, порождая переходный квант
$\Delta_4 - \Delta_3 = 15 - 4 = 11$.
\item[(iii)] \textbf{Защита.} Граница Хэмминга $2^r \ge 4 + r + 1$ фиксирует
$r = 3$; семь ненулевых синдромов изоморфны плоскости Фано PG$(2,2)$ и образуют
замкнутую коммутативную идемпотентную квазигруппу Штейнера.
\item[(iv)] \textbf{Период и сброс.} Тактовый цикл $T_0 = 12 = 3 \times 4 = |A_4|$,
модуль сброса $13$ ($\gcd(13,12) = 1$), коридор разрядности $d_{\mathrm{bit}} = 7 \pm 1$;
расширение до $n = 5$ блокировано: $\subfact{5} = 44 = 4 \cdot 11$ не является
степенью простого.
\item[(v)] \textbf{Стрела времени.} Каналы удержания инварианта
\[
W_1(3 \to 4) = \frac{13}{16}, \qquad
W_1(4 \to 3) = \frac{13}{15}
\]
необратимы: принцип детального баланса нарушен с точным коэффициентом
\[
\chi = \frac{W_1(4 \to 3)}{W_1(3 \to 4)} = \frac{16}{15} \ne 1 .
\]
Стационарное распределение двухтактной цепи равно
$\left(\frac{16}{31},\,\frac{15}{31}\right)$, а квант асимметрии потерь инварианта
\[
\frac{3}{16} - \frac{2}{15} = \frac{13}{240} > 0
\]
ориентирует время: расширение $3 \to 4$ теряет инвариант строго чаще стягивания
$4 \to 3$, то есть стрела направлена в сторону роста потери инвариантной
информации.
\end{enumerate}
\end{theorem}

\section{Леммы о структуре провалов каналов}

Взаимодействие каналов уточняется следующей структурной леммой.

\begin{lemma}[Провалы прямого канала]\label{lm:fwd}
Пусть $\pi \in M_1(3)$, $\pi^{(4)} = \pi \cup \{4 \mapsto 4\}$, и вставка
определена как $\sigma = (i\;4) \circ \pi^{(4)}$, $i \in \{1,2,3,4\}$. Тогда
$\sigma \notin M_1(4)$ тогда и только тогда, когда $\pi$ --- транспозиция и
транспозиция $(i\;4)$ не пересекается с носителем её нетривиального цикла.
В частности, число провалов равно $3$ (по одному на каждую транспозицию
$S_3$), и $W_1(3 \to 4) = 1 - \frac{3}{16} = \frac{13}{16}$.
\end{lemma}

\begin{proof}
Произведение двух транспозиций $(i\;4)\tau$ не имеет неподвижных точек тогда и
только тогда, когда они не пересекаются; в любом другом случае неподвижная
точка сохраняется (в т.~ч. при $\pi = \mathrm{id}$ и при пересекающихся
транспозициях, где цикл склеивается в 3-цикл с неподвижной точкой). Для каждой
из трёх транспозиций $\tau \subset S_3$ существует ровно одно $i$ с
$\{i,4\} \cap \mathrm{supp}(\tau) = \emptyset$. Итого провалов $3$ из $16$.
\end{proof}

\begin{lemma}[Провалы обратного канала]\label{lm:bwd}
Ограничение $\rho(\sigma) \in \S_3$ (удаление элемента $4$ из цикловой
структуры) теряет неподвижную точку тогда и только тогда, когда $\sigma|_{\{1,2,3\}}$
--- 3-цикл и $\sigma(4) = 4$. Таких $\sigma \in M_1(4)$ ровно две
($3$-циклы $(123)$ и $(132)$), поэтому
$W_1(4 \to 3) = 1 - \frac{2}{15} = \frac{13}{15}$.
\end{lemma}

\begin{proof}
Если $\pi = \rho(\sigma)$ не имеет неподвижных точек в $\S_3$, то $\pi$ ---
3-цикл; тогда все значения $1,2,3$ заняты этим циклом и биективность вынуждает
$\sigma(4) = 4$. Обратно, для $\sigma = \pi \cup \{4 \mapsto 4\}$ с 3-циклом
$\pi$ ограничение теряет неподвижную точку. Во всех прочих случаях удаление
$4$ из цикла склеивает соседей, и хотя бы одна неподвижная точка сохраняется.
\end{proof}

\begin{lemma}[Нарушение детального баланса]\label{lm:db}
Для двухтактной цепи с интенсивностями $r_{34} = \frac{13}{16}$,
$r_{43} = \frac{13}{15}$: (а) $\chi = r_{43}/r_{34} = \frac{16}{15}$;
(б) стационарное распределение $w = \left(\frac{16}{31},\,\frac{15}{31}\right)$,
то есть $w_3/w_4 = \frac{16}{15}$; (в) при равных весах ансамблей чистый поток
$J = \frac{1}{2}\left(\frac{13}{16} - \frac{13}{15}\right) = -\frac{13}{240} \ne 0$.
\end{lemma}

\begin{proof}
Для двух состояний $w_3/w_4 = r_{43}/r_{34} = \frac{13/15}{13/16} = \frac{16}{15}$;
нормировка $w_3 + w_4 = 1$ даёт $\frac{16}{31}, \frac{15}{31}$. Остальное --
прямое вычисление.
\end{proof}

\begin{remark}[Структурная роль числа 13]
Число $13$ возникает в теории четырежды, и теперь все четыре возникновения
имеют доказательную, а не декларативную природу:
\[
13 = 16 - \#\{\text{транспозиций в } S_3\} = 15 - \#\{\text{3-циклов в } S_3\}
\]
--- счётчик сохранённых состояний в обоих каналах; модуль сброса PG$(2,3)$;
число точек проективной прямой $\mathbb{P}^1(\mathbb{F}_{11})$ плюс единица;
числитель кванта асимметрии $\frac{13}{240}$.
\end{remark}

\begin{remark}[Метрический квант]
Суммирование информационной метрики Фишера---Рао---Ченцова по двунаправленному
циклу даёт $\Delta\tau_{\mathrm{metric}} = \ln\!\left(\frac{7}{4}\left(\frac{15}{16}\right)^{11/6}\right)
\approx 0.441295$ нат --- численно подтверждено.
\end{remark}

\begin{remark}[Границы теоремы]
Пункты (i)--(v) --- математические утверждения о конечных объектах
($\S_n$, коды, плоскости Фано, марковские цепи) и доказаны полностью.
Идентификация полученной структуры с физическим временем --- содержательная
гипотеза теории; теорема фиксирует её математическое ядро, свободное от
свободных параметров.
\end{remark}

\end{document}
